%%%%%%%%%%%%%%%%%%%%%%%%%%%%%%%%%%%%%%%%%%%%%%%%%%%%%%%%%%%%%%%%%%%%%%
\section{Contexte et motivation}
%%%%%%%%%%%%%%%%%%%%%%%%%%%%%%%%%%%%%%%%%%%%%%%%%%%%%%%%%%%%%%%%%%%%%%

Depuis plusieurs années, les progrès réalisés en intelligence artificielle
ont profondément transformé le traitement automatique de la parole et des
langues. Des modèles tels que Whisper, wav2vec~2.0, HuBERT ou encore MMS
(Massively Multilingual Speech) ont démontré qu'il était désormais possible
d'apprendre des représentations linguistiques riches directement à partir
de grands volumes de données audio.

Cependant, ces avancées concernent principalement des langues disposant de
ressources importantes : corpus annotés, dictionnaires, systèmes
d'écriture standardisés et importantes quantités de données numériques.

À l'inverse, de nombreuses langues africaines restent encore très peu
représentées dans les technologies actuelles. Elles sont souvent qualifiées
de \emph{langues peu dotées} (\emph{low-resource languages}), car elles ne
disposent que de très peu de ressources exploitables pour entraîner des
modèles d'intelligence artificielle.

Le soussou, parlé principalement en Guinée, constitue un exemple
particulièrement intéressant. Malgré un nombre important de locuteurs,
cette langue demeure largement sous-représentée dans les grands corpus
internationaux utilisés aujourd'hui pour le développement des modèles de
traitement de la parole.

\vspace{0.5cm}

Cette situation soulève plusieurs difficultés.

Les systèmes modernes de reconnaissance vocale ne prennent généralement
pas en charge le soussou.

Les modèles de traduction automatique sont très limités.

Les assistants vocaux actuels sont incapables de comprendre ou de produire
cette langue.

Enfin, les ressources linguistiques disponibles sont insuffisantes pour
développer rapidement de nouvelles applications.

\vspace{0.5cm}

Pourtant, les récents progrès des modèles auto-supervisés montrent qu'il
n'est plus nécessaire de disposer d'énormes quantités de données annotées
pour apprendre des représentations utiles de la parole.

Cette évolution ouvre de nouvelles perspectives pour les langues peu
dotées et conduit naturellement à la question suivante :

\begin{center}
\fbox{
\parbox{0.9\textwidth}{
\centering
\textbf{
Une intelligence artificielle peut-elle apprendre progressivement une
langue essentiellement orale à partir de données de parole collectées
auprès de locuteurs natifs, puis utiliser cette connaissance pour assister
à la création d'outils linguistiques adaptés à cette langue ?
}
}
}
\end{center}

Cette interrogation constitue le point de départ du projet
\textbf{ALORA}.

\vspace{0.5cm}

Le choix du soussou comme premier cas d'étude ne repose pas uniquement sur
le fait qu'il s'agisse de ma langue maternelle. Il répond également à une
motivation scientifique.

En effet, travailler sur une langue peu dotée permet d'étudier les limites
des méthodes actuelles d'apprentissage automatique lorsque les ressources
sont rares, tout en développant des méthodologies susceptibles d'être
réutilisées pour d'autres langues présentant des caractéristiques
similaires.

Ainsi, le soussou constitue un terrain d'expérimentation pertinent pour
explorer de nouvelles approches en traitement automatique de la parole et
en linguistique computationnelle.

\vspace{0.5cm}

Ce projet ne cherche donc pas uniquement à développer une application
destinée au soussou.

L'ambition est plus générale : proposer une méthodologie permettant de
faciliter la création de ressources numériques pour des langues orales peu
dotées, en combinant collecte collaborative de données, intelligence
artificielle, traitement de la parole et apprentissage automatique.

Cette réflexion s'inscrit dans une perspective de recherche à long terme,
dont la première étape consiste à construire des bases méthodologiques
solides avant le développement des modèles d'intelligence artificielle.